\uuid{FNYm}
\titre{Théorème de Bézout}
\auteur{bodin ; croizat ; sacre}
\organisation{exo7}
\licence{CC BY-NC-SA 4.0}
\datecreate{2022-03-08}
\niveau{L1}
\module{Algèbre}
\chapitre{Arithmétique dans Z}
\sousChapitre{Nombres premiers, nombres premiers entre eux}
\difficulte{2}
\theme{Théorème de Bézout}
\exo7qcmid{hilisit.arithmetique.019}

\contenu{
\begin{question}
\qclassification{103.04}
    Soient deux entiers positifs $a,b$ et $d=\pgcd(a,b)$.
    Quelles sont les affirmations vraies ?
    \begin{answers}
        \bad{Il existe $u,v\in\Zz$ uniques tels que $au+bv=d$.}
        \good{Il existe $u,v\in\Zz$ tels que $au+bv=d$.}
        \bad{Il existe $u,v\in\Nn$ uniques tels que $au+bv=d$.}
        \bad{Il existe $u,v\in\Nn$ tels que $au+bv=d$.}
    \end{answers}
    \begin{explanations}
     Il existe $u,v\in\Zz$ tels que $au+bv=d$. $u$ et $v$ ne sont pas uniques. Comme $a,b,d>0$, $d \le a$ et $d \le b$ alors soit $u$ soit $v$ sera négatif.
    \end{explanations}
\end{question}
}
